\section{Discussion}

The paper places a narrow causal object on a complete mathematical and inferential footing. The intervention is an operator-module replacement inside a declared linear Gaussian SPDE class. The target is a counterfactual field-functional distribution. Identification is target specific. The inferential limit is mixed rate. Numerical approximation must be controlled at the target family's root-information scale.

Three distinctions are particularly important.

First, semantics and identification are separate. \Cref{sec:semantics} provides a coherent shared-noise intervention even before one asks whether the resulting law is recoverable from factual data; \Cref{sec:identification,sec:nonid} then characterise the positive and negative identification results.

Second, strong field approximation and Gaussian law approximation need not have the same rate. Under non-degeneracy, the variance tail enters normal CDF, quantile, and optimal Wasserstein comparisons linearly. Using the common-noise strong rate for every law target can be unnecessarily conservative.

Third, the phrase ``causal SPDE'' should not be read as unrestricted causal discovery. The intervention and invariant modules must be supplied. The feedback-gauge and open-loop constructions show that passive equality of factual laws cannot decide which part of a drift is the module that should be replaced.

Several extensions remain. Non-diagonal unbounded perturbations, including spatially varying divergence-form coefficients, require domain-sensitive perturbation theory. Semilinear equations require differentiability of nonlinear solution maps. Non-Gaussian noise calls for characteristic-function, Wasserstein, or quantile-specific methods rather than Gaussian mean--variance reduction. Physical-space observations require a separate statistical experiment and may produce different rates.

The present results are nevertheless sufficient for a coherent paper-level conclusion:
\begin{equation*}
\begin{aligned}
&\text{specified structural intervention}
+\text{identified factual experiment}\\
&\quad+\text{controlled numerical approximation}\\
&\hspace{2.5cm}\Longrightarrow
\text{uniformly inferable counterfactual laws}.
\end{aligned}
\end{equation*}
The impossibility results define the boundary of that conclusion rather than an exception to it.
